\chapter{Năm Phút Mở Đường Cho Cả Tiết}
\chapterintro{Khởi động đúng để cả tiết đỡ trôi}{Năm phút đầu không chỉ để gọi lớp chú ý, mà còn để đặt đường ray cho toàn bộ tiết học phía sau.}

\section{Nêu mục tiêu bằng một câu}
\begin{khoidongbox}{Hoạt động khởi động}
Mở tiết bằng đúng một câu nói rõ hôm nay lớp cần đi tới đâu.
\end{khoidongbox}
\begin{visaohieuqua}
Khi biết đích đến, học sinh bớt cảm giác bị kéo qua một tiết học mơ hồ. Mục tiêu rõ giúp lớp hợp tác hơn.
\end{visaohieuqua}
\begin{cachlam5phut}
Mục tiêu phải ngắn, nói theo ngôn ngữ của lớp, và gắn với việc các em sẽ làm trong tiết.
\end{cachlam5phut}
\ngancach

\section{Cho lớp chọn đường vào bài}
\begin{khoidongbox}{Hoạt động khởi động}
Đưa hai đường vào bài như: ví dụ trước hay khái niệm trước, câu hỏi trước hay tình huống trước.
\end{khoidongbox}
\begin{visaohieuqua}
Không phải bài nào cũng cần một lối mở cố định. Khi được chọn đường vào, học sinh cảm thấy mình có vai trò ngay từ đầu.
\end{visaohieuqua}
\begin{cachlam5phut}
Cho lớp giơ tay chọn đa số, rồi đi theo hướng đó nhưng vẫn giữ khung nội dung chính.
\end{cachlam5phut}
\ngancach

\section{Mời một em mở bằng ngôn ngữ của lớp}
\begin{khoidongbox}{Hoạt động khởi động}
Mời một học sinh tóm nội dung sắp học bằng cách nói của chính các em, thật ngắn và tự nhiên.
\end{khoidongbox}
\begin{visaohieuqua}
Ngôn ngữ của học sinh thường làm cả lớp thấy bài gần mình hơn. Nó cũng giúp thầy cô đo trước mức hiểu và mức hứng thú.
\end{visaohieuqua}
\begin{cachlam5phut}
Chọn một em bất kỳ hoặc một em đã chuẩn bị. Nghe xong thì chốt lại bằng ngôn ngữ chính xác hơn nếu cần.
\end{cachlam5phut}
\ngancach

\section{Câu hỏi chờ cả tiết}
\begin{khoidongbox}{Hoạt động khởi động}
Đặt ra một câu hỏi sẽ được giữ lại suốt tiết và hứa cuối giờ quay về trả lời trọn vẹn.
\end{khoidongbox}
\begin{visaohieuqua}
Một câu hỏi treo khiến lớp có lý do để đi hết tiết. Các em học trong trạng thái chờ kiểm chứng nên đỡ trôi hơn.
\end{visaohieuqua}
\begin{cachlam5phut}
Câu hỏi phải vừa sức nhưng đủ gợi. Viết nó lên bảng và nhắc lại ở giữa hoặc cuối tiết.
\end{cachlam5phut}
\ngancach

\section{Chốt mốc đầu tiết}
\begin{khoidongbox}{Hoạt động khởi động}
Trước khi vào nội dung chính, chốt một mốc đầu tiết: điều lớp phải nhớ, phải làm, hoặc phải tránh nhầm.
\end{khoidongbox}
\begin{visaohieuqua}
Nhiều tiết học trôi vì đầu tiết không có mốc bám. Một mốc rõ giúp học sinh biết mình đang đi trên đường nào.
\end{visaohieuqua}
\begin{cachlam5phut}
Nói mốc đó bằng một câu ngắn, có thể viết lên bảng. Đó sẽ là điểm quay lại nếu lớp bị trôi ở giữa tiết.
\end{cachlam5phut}